\clearpage\thispagestyle{empty}
\pdfbookmark[0]{Quellen zur Besprechung von Hardings Himmelskarten}{harding1817-sources}
\begin{center}\Large\bfseries Quellen zur Besprechung von Hardings Himmelskarten\end{center}

Die Übertragung folgt Carl Friedrich Gauß’ Besprechung in den \emph{Werken}, Bd.~VI (Göttingen, 1874), S.~593--595. Der Erstabdruck erschien in den \emph{Göttingischen gelehrten Anzeigen}, Stück~5, 9.~Januar 1817, S.~41--45. Die Besprechung verweist auf die Anzeige in Stück~158 derselben Zeitschrift von 1812.

Die Werkausgabe hat \emph{Athas}, der Erstabdruck \emph{Atlas}. In der Werkausgabe steht nach \emph{ist} ein Komma, im Erstabdruck ein Punkt. Die Übertragung übernimmt beide Lesarten der Werkausgabe. Die Schreibweise \emph{3osten} gibt das runde, o-förmige Zeichen der Werkausgabe wieder. Die historische Beschreibung der Projektion folgt dem Wortlaut der Besprechung. Das historische \emph{unentgeldlich} erscheint in moderner Schreibweise als \emph{unentgeltlich}.

Gauß’ Besprechung kündigt insgesamt 28 Blätter an. Das Titelblatt von Karl Ludwig Hardings \emph{Atlas novus coelestis} (Göttingen, 1822) nennt 27 Blätter. Die Karten werden in der Besprechung beschrieben.

KI-Übertragung aus dem historischen Deutsch in modernes Deutsch. KI-erstellter Satz.

\href{https://gdz.sub.uni-goettingen.de/id/PPN236007467}{Werkausgabe}. Erstdruck, S.~\href{https://images.p.sub.uni-goettingen.de/iiif/image/gdz:PPN319721507_1817:00000043/full/1650,/0/default.jpg}{41}, \href{https://images.p.sub.uni-goettingen.de/iiif/image/gdz:PPN319721507_1817:00000044/full/1650,/0/default.jpg}{42}, \href{https://images.p.sub.uni-goettingen.de/iiif/image/gdz:PPN319721507_1817:00000045/full/1650,/0/default.jpg}{43}, \href{https://images.p.sub.uni-goettingen.de/iiif/image/gdz:PPN319721507_1817:00000046/full/1650,/0/default.jpg}{44}, \href{https://images.p.sub.uni-goettingen.de/iiif/image/gdz:PPN319721507_1817:00000047/full/1650,/0/default.jpg}{45}. \href{https://images.p.sub.uni-goettingen.de/iiif/image/gdz:PPN319721507_1812:00001613/full/1650,/0/default.jpg}{Anzeige von 1812}. \href{https://www.davidrumsey.com/luna/servlet/detail/RUMSEY~8~1~372238~90138977:Title--Atlas-novus-coelestis-XXVII-}{Titelblatt des Atlas von 1822}.

\clearpage
